\section{Failure of Local Rules}

We show that the transport invariant cannot be recovered from purely
local combinatorial data on $\Gfifteen$.

\subsection{Local indistinguishability}

The graph $\Gfifteen$ is $4$-regular and edge-transitive. In particular,
for any fixed radius $r$, all edges have isomorphic $r$-neighborhoods.

Thus no invariant depending only on finite-radius neighborhoods and
respecting graph automorphisms can distinguish one edge from another.

\subsection{Global variability of the cocycle}

Despite this local symmetry, the Thalean cocycle
\[
\varepsilon : E(\Gfifteen) \to \mathbb{Z}_2
\]
takes both values $0$ and $1$.

Therefore the transport invariant distinguishes edges which are locally
indistinguishable.

\subsection{Failure of local determination}

\begin{theorem}
The Thalean cocycle is not determined by any rule depending only on
finite-radius neighborhoods in $\Gfifteen$ and invariant under graph
automorphisms.
\end{theorem}

\begin{proof}
Any such rule assigns values to edges based only on their
$r$-neighborhoods up to isomorphism and is invariant under automorphisms.

Since $\Gfifteen$ is edge-transitive, all edges have isomorphic
$r$-neighborhoods. Hence any such rule must assign a constant value to
all edges.

However, the cocycle $\varepsilon$ takes both values $0$ and $1$ (see
Section~13). Therefore it cannot arise from any such local rule.
\end{proof}

\subsection{Geometric origin of nonlocality}

The failure of locality is not accidental. The cocycle is defined by
lifting edges to the double cover $G_{30}\to G_{15}$ and then resolving
those lifted edges in the canonical dodecahedral model.

In that model, each edge of $\Gfifteen$ corresponds to a pair of
relations between dodecahedral edges, and the cocycle records whether
these relations are both local or both nonlocal.

This information is invisible in $\Gfifteen$ itself and only becomes
accessible after passing to the lifted structure and its geometric
realization.

\subsection{Cohomological interpretation}

The obstruction to locality is cohomological. The cocycle $\varepsilon$
represents a nontrivial class in $H^1(\Gfifteen,\ZZtwo)$ and cannot be
expressed as the coboundary of a vertex potential.

Equivalently, no vertex switching can trivialize the transport invariant.

\subsection{Consequences}

The transport invariant encodes genuinely global information: it is
determined by lifted transport and geometric incidence in the
dodecahedron, and cannot be reconstructed from local combinatorial data
on the core graph alone.

Thus the Petersen line-graph core carries a hidden global structure
revealed only through the Thalean lift.
